\section{Series reduction and boundary cases}\label{sec:examples}

\subsection{Chains and rays}
The weight in \eqref{eq:potential} satisfies the exact identity
\begin{equation}\label{eq:series}
 \frac{1-r_e^2}{R_w^2}=\frac1{R_w^2}-\frac1{R_v^2},
 \qquad e=(v,w).
\end{equation}
Along a finite chain whose internal vertices have one child, a conserved
flow has constant value. Its potential contribution therefore telescopes
to the contribution of one edge with channel parameter equal to the
product of the parameters on that chain. The channels themselves obey
$P_sP_t=P_{st}$. Thus replacing such a chain by its composite channel
preserves the capacity, as well as the experiment between its endpoints.
For a constant parameter $0<\lambda<1$, the capacity here equals
$(1-\lambda^4)^{-1/2}$ times the capacity normalized in \cite{Capacity}.

On a single infinite ray, put $P_n=\prod_{j=1}^n\lambda_j$ and
$P_\infty=\lim_nP_n$. The unit flow has potential
\[
 \lim_{n\to\infty}\bigl(P_n^{-4}-1\bigr).
\]
Consequently,
\begin{equation}\label{eq:ray-capacity}
 C(T)=
 \begin{cases}
  0,&P_\infty=0,\\
  (P_\infty^{-4}-1)^{-1/2},&0<P_\infty<1,\\
  \infty,&P_\infty=1.
 \end{cases}
\end{equation}
The depth-$n$ experiment is exactly the channel $P_{P_n}$, so
\[
 m_n=3P_n^2,\qquad a_n=\frac34P_n.
\]
This verifies the criterion on a ray directly, without
Assumption~\ref{ass:input}.

For example, let $\lambda_j=1-(j+1)^{-2}$ for $j\geq1$.
Then
\[
 P_n=\frac{n+2}{2(n+1)}\longrightarrow\frac12,
 \qquad C(T)=\frac1{\sqrt{15}},\qquad a_\infty=\frac38.
\]
In contrast, simply replacing $\lambda^{2|v|}$ by $R_v$ in
the unweighted fixed-channel potential gives the divergent series
$\sum_nP_n^{-4}$, and hence zero capacity on this ray. The factor
$1-\lambda_e^4$ is therefore needed when parameters approach one.
A perfectly transmitting ray has zero potential for every finite flow,
so $C(T)=\infty$ and its spin observations reveal the root exactly.

\subsection{Zero channels}
\begin{corollary}\label{cor:zero}
Assume Assumption~\ref{ass:input}, and allow
$0\leq\lambda_e\leq1$. Let $T_+$ be the connected component of
the root after deleting all zero-channel edges. Define capacity using
\eqref{eq:potential}--\eqref{eq:capacity} on $T_+$. Then
Theorem~\ref{thm:main}, including its quantitative bounds, remains valid
with $C(T_+)$ in place of $C(T)$.
\end{corollary}
\begin{proof}
Realize each channel by copying the parent with probability
$\lambda_e$ and otherwise drawing a fresh uniform spin, using
independent edge randomness. Across a zero edge the child is a fresh
uniform spin independent of its parent. Group removed vertices by the
first zero edge on their path from the root. Each resulting subtree
depends only on its fresh initial spin and subsequent independent
randomness. Jointly, all their observations are independent of both
the root and all observations in $T_+$. Deleting them leaves the
root posterior at every depth unchanged. Apply the theorem to the
positive-channel tree $T_+$. If it is finite, distant levels are empty
and conservation forces every flow to vanish.
\end{proof}

\subsection{The role of local finiteness}
Local finiteness cannot be omitted without changing the statement.
Give the root, for each $k\geq1$, a disjoint child chain of length $k$,
and take every channel to be perfect. There is no infinite ray.
Conservation forces the flow on each finite chain to vanish, working
backward from its leaf. Even with countably additive conservation at
the root, every flow has mass zero. Yet each level $n$ contains a spin
equal to the root, from any chain of length at least $n$. Thus
$m_n=3$ and $a_n=3/4$ for every $n$, although the capacity is zero.
